\subsubsection{Arzela-Ascoli 定理}\label{subsec:arzela-ascoli-theorem}

Arzela-Ascoli 定理指出, 紧集上满足某些条件的函数序列具有一致收敛子列. 这类结果在存在性证明中常常有用; 需要证明存在的对象就是收敛子列的极限. 为陈述该定理, 需要以下定义.

\begin{Definition}
\label{def:equicontinuity}
设 $f_m$ 为定义在 $\mathbb{R}^n$ 的子集 $D$ 上的实值函数序列. 设 $\bm{x}\in D$. 若对每个 $\epsilon>0$, 都存在不依赖于 $m$ 的 $\delta>0$, 使对满足 $|\bm{y}-\bm{x}|<\delta$ 的 $\bm{y}\in D$ 都有 $|f_m(\bm{y})-f_m(\bm{x})|<\epsilon$, 则称该序列在 $\bm{x}$ 处\textbf{等度连续}.
\end{Definition}

若序列 $f_m$ 在紧集 $S$ 的每一点处都等度连续, 则它一致等度连续, 即上述定义中的 $\delta$ 可以不依赖于 $\bm{x}\in S$ 来选取 (参见 Problem~\ref{ex:compact-equicontinuity-uniform}; 不要求 $D=S$). 注意, 若在 $\bm{x}$ 的某个邻域内导数具有不依赖于 $m$ 的界, 则函数序列在 $\bm{x}$ 处等度连续.

\begin{Theorem}{Arzela-Ascoli}
\label{thm:arzela-ascoli}
设 $f_m$ 为定义在 $\mathbb{R}^n$ 的紧子集 $S$ 上的实值函数序列. 假设存在常数 $M$, 使对每个 $m\in\mathbb{N}$ 和每个 $\bm{x}\in S$ 都有 $|f_m(\bm{x})|\leq M$. 再假设序列 $f_m$ 在 $S$ 的每一点处等度连续. 那么存在一个在 $S$ 上一致收敛的子列.
\end{Theorem}

指出, 经过相当明显的修改, 定理和证明可以推广到 $S$ 为抽象拓扑空间中的紧集, 且 $f_m$ 的值位于完备度量空间中的情形.

\begin{Proof}
设 $\bm{x}_i$, $i\in\mathbb{N}$ 为在 $S$ 中稠密的点列. 序列 $f_m(\bm{x}_1)$ 有界, 因而有收敛子列. 即可选择序列 $m_{1j}$, 使 $f_{m_{1j}}(\bm{x}_1)$ 当 $j\to\infty$ 时收敛. 类似地, 可选取序列 $m_{1j}$ 的子列 $m_{2j}$, 使 $f_{m_{2j}}(\bm{x}_2)$ 收敛. 由于 $m_{2j}$ 是 $m_{1j}$ 的子列, $f_{m_{2j}}(\bm{x}_1)$ 也收敛. 接着选取序列 $m_{2j}$ 的子列 $m_{3j}$, 使 $f_{m_{3j}}$ 在 $\bm{x}_3$ 处也收敛. 如此无限继续. 最后考虑"对角"序列 $f_{m_{jj}}$. 除去前 $i-1$ 项, $m_{jj}$ 是 $m_{ij}$ 的子列; 因此 $f_{m_{jj}}(\bm{x}_i)$ 对每个 $i\in\mathbb{N}$ 都收敛. 为简化记号, 下文令 $g_j=f_{m_{jj}}$.

为完成证明, 说明序列 $g_m$ 是一致 Cauchy 的. 给定 $\epsilon>0$. 由于 $g_m$ 是 $f_m$ 的子列, 它们在 $S$ 上一致等度连续; 因而存在 $\delta>0$, 使只要 $|\bm{y}-\bm{x}|<\delta$, 就有 $|g_m(\bm{y})-g_m(\bm{x})|<\epsilon/3$. 由于 $S$ 紧, 存在 $K\in\mathbb{N}$, 使对每个 $\bm{x}\in S$, 都存在 $i\in\{1,\ldots,K\}$ 满足 $|\bm{x}_i-\bm{x}|<\delta$. 现在将 $N$ 取得足够大, 使对 $m,k>N$ 及每个 $i\in\{1,\ldots,K\}$ 都有 $|g_m(\bm{x}_i)-g_k(\bm{x}_i)|<\epsilon/3$. 对 $m,k>N$ 及任意 $\bm{x}\in S$, 有
\begin{equation}\label{eq:arzela-ascoli-uniform-cauchy-estimate}
\begin{aligned}
|g_m(\bm{x})-g_k(\bm{x})|
&\leq|g_m(\bm{x})-g_m(\bm{x}_i)|\\
&\quad+|g_m(\bm{x}_i)-g_k(\bm{x}_i)|+|g_k(\bm{x}_i)-g_k(\bm{x})|\\
&<\epsilon,
\end{aligned}
\end{equation}
其中 $i\in\{1,\ldots,K\}$ 为某个指标.
\end{Proof}

下文需要对调和函数序列应用 Arzela-Ascoli 定理. 对这一特殊情形, 有以下结果.

\begin{Theorem}
\label{thm:harmonic-sequence-compactness}
设 $\Omega$ 为 $\mathbb{R}^n$ 中的区域. 设 $f_m$ 为 $\Omega$ 上一致有界的调和函数序列, 即对每个 $\bm{x}\in\Omega$ 和 $m\in\mathbb{N}$ 都有 $|f_m(\bm{x})|\leq M$. 那么 $f_m$ 有一个子列, 收敛到 $\Omega$ 上的调和函数, 并在 $\Omega$ 的紧子集上一致收敛.
\end{Theorem}

\begin{Proof}
令 $\Omega_k=\{\bm{x}\in\Omega\mid|\bm{x}|\leq k,\ \operatorname{dist}(\bm{x},\partial\Omega)\geq1/k\}$. 那么 $\Omega_k$ 是紧的, 且 $\Omega=\bigcup_{k=1}^\infty\Omega_k$. 若 $\bm{x}$ 位于 $\Omega_k$ 中, 则在 $\bm{x}$ 的一个邻域内, 可用以 $\bm{x}$ 为球心, 半径为 $1/2k$ 的球, 通过 Poisson 公式表示 $f_m$. 这给出 $f_m$ 在 $\Omega_k$ 上的导数的一致估计. 特别地, $f_m$ 在 $\Omega_k$ 上等度连续. 由 Arzela-Ascoli 定理, 可在每个 $\Omega_k$ 上提取一致收敛子列, 再使用与上述证明类似的对角论证, 得到在整个 $\Omega$ 上收敛, 且在每个 $\Omega_k$ 上一致收敛的子列.

剩下只需证明这一子列的极限调和. 为此, 只需注意, 调和函数的特征性质是在每个球上都满足 Poisson 公式. 若一致收敛序列中的所有函数都满足 Poisson 公式, 则极限也满足该公式.
\end{Proof}
